\documentclass[aspectratio=169,11pt]{beamer}
\usepackage{../../../shared/theme/esmad_beamer_theme}
\usepackage{amsmath,amssymb}
\usepackage{booktabs}
\usepackage{array}

\title[Panel Data Econometrics]{Panel Data, Fixed Effects, and Unobserved Heterogeneity}
\subtitle{Identification from Repeated Observations}
\date{\today}

\newcommand{\E}{\mathbb{E}}
\newcommand{\Var}{\mathrm{Var}}
\newcommand{\Cov}{\mathrm{Cov}}
\newcommand{\plim}{\operatorname{plim}}

\begin{document}

\begin{frame}
\titlepage
\end{frame}

\begin{frame}{Learning Goals}
\begin{itemize}
\item explain what panel data add beyond repeated cross-sections;
\item derive the within transformation and interpret fixed-effects estimators;
\item distinguish pooled OLS, fixed effects, random effects, and two-way fixed effects;
\item state the exogeneity assumptions required for causal interpretation;
\item choose uncertainty estimates that respect panel dependence;
\item recognize time-varying confounding, dynamic feedback, and other failures that fixed effects do not solve.
\end{itemize}
\end{frame}

\begin{frame}{Outline}
\tableofcontents
\end{frame}

\section{Why Panel Data?}

\begin{frame}{The Panel Data Structure}
A balanced panel observes units \(i=1,\ldots,N\) over periods \(t=1,\ldots,T\):
\[
\{Y_{it},D_{it},X_{it}\}.
\]

\begin{columns}[T]
\begin{column}{0.48\textwidth}
\begin{block}{Cross-sectional variation}
Why do units differ from one another?
\end{block}
\end{column}
\begin{column}{0.48\textwidth}
\begin{block}{Within-unit variation}
How does the same unit change over time?
\end{block}
\end{column}
\end{columns}

\vspace{0.4cm}
Panel methods often identify coefficients primarily from within-unit changes after removing persistent heterogeneity.
\end{frame}

\begin{frame}{A Canonical Panel Model}
Consider
\[
Y_{it}=\beta D_{it}+X_{it}'\gamma+\alpha_i+u_{it},
\]
where:
\begin{itemize}
\item \(\alpha_i\) is time-invariant unit heterogeneity;
\item \(u_{it}\) is the remaining time-varying disturbance;
\item \(D_{it}\) is the treatment or exposure of interest.
\end{itemize}

\begin{alertblock}{Central identification issue}
If \(\alpha_i\) is correlated with \(D_{it}\), pooled OLS generally mixes treatment variation with persistent unit differences.
\end{alertblock}
\end{frame}

\section{Pooled OLS and Fixed Effects}

\begin{frame}{Why Pooled OLS Can Fail}
Pooled OLS ignores unit structure:
\[
Y_{it}=\beta D_{it}+X_{it}'\gamma+\varepsilon_{it},
\qquad
\varepsilon_{it}=\alpha_i+u_{it}.
\]

If
\[
\Cov(D_{it},\alpha_i)\neq 0,
\]
then
\[
\Cov(D_{it},\varepsilon_{it})\neq 0
\]
even when the idiosyncratic disturbance is otherwise well behaved.

\begin{block}{Example}
More productive firms may adopt a technology earlier. A cross-sectional regression then confounds technology adoption with persistent productivity differences.
\end{block}
\end{frame}

\begin{frame}{The Within Transformation}
Start from
\[
Y_{it}=\beta D_{it}+X_{it}'\gamma+\alpha_i+u_{it}.
\]

Take the unit-specific time mean:
\[
\bar Y_i=\beta \bar D_i+\bar X_i'\gamma+\alpha_i+\bar u_i.
\]

Subtract:
\[
Y_{it}-\bar Y_i
=
\beta(D_{it}-\bar D_i)
+
(X_{it}-\bar X_i)'\gamma
+
(u_{it}-\bar u_i).
\]

\begin{alertblock}{Result}
The time-invariant \(\alpha_i\) disappears exactly.
\end{alertblock}
\end{frame}

\begin{frame}{What Variation Identifies the Fixed-Effects Coefficient?}
The FE estimator uses deviations from each unit's own mean:
\[
\widetilde D_{it}=D_{it}-\bar D_i.
\]

Thus \(\beta\) is learned from units whose treatment changes over time.

\begin{itemize}
\item units with no within-unit treatment variation contribute no identifying variation for \(\beta\);
\item time-invariant regressors are absorbed by unit fixed effects;
\item interpretation is fundamentally within-unit, not cross-sectional.
\end{itemize}

\begin{block}{Practical question}
Is the within-unit variation in treatment substantively meaningful, or mostly noise and measurement error?
\end{block}
\end{frame}

\section{Exogeneity Assumptions}

\begin{frame}{Strict Exogeneity}
A standard FE condition is
\[
\E[u_{it}\mid D_{i1},\ldots,D_{iT},X_{i1},\ldots,X_{iT},\alpha_i]=0
\]
for every \(t\).

This rules out correlation between today's disturbance and treatment in any period.

\begin{alertblock}{Implication}
Future treatment cannot respond to today's unexpected outcome shock under strict exogeneity.
\end{alertblock}

That can be implausible when treatment is adjusted dynamically in response to past outcomes.
\end{frame}

\begin{frame}{Contemporaneous Exogeneity Is Weaker}
A weaker condition is
\[
\E[u_{it}\mid D_{it},X_{it},\alpha_i]=0.
\]

This permits some temporal feedback but is generally insufficient for the standard static FE estimator when treatment decisions respond to lagged shocks.

\begin{block}{Example}
A hospital increases staffing next month because this month's unobserved demand shock was unusually large.
\end{block}

The panel structure does not eliminate this feedback mechanism.
\end{frame}

\section{Two-Way Fixed Effects}

\begin{frame}{Adding Time Fixed Effects}
A two-way fixed-effects model is
\[
Y_{it}=\beta D_{it}+X_{it}'\gamma+\alpha_i+\lambda_t+u_{it}.
\]

Here:
\begin{itemize}
\item \(\alpha_i\) absorbs stable unit-specific factors;
\item \(\lambda_t\) absorbs common shocks at time \(t\).
\end{itemize}

\begin{block}{Examples of common time shocks}
Inflation, national policy changes, seasonality, macroeconomic cycles, or platform-wide product changes.
\end{block}
\end{frame}

\begin{frame}{What Two-Way Fixed Effects Still Do Not Remove}
TWFE does not automatically remove:
\begin{itemize}
\item unit-specific time-varying confounders;
\item heterogeneous trends across units;
\item anticipation before treatment;
\item spillovers across units;
\item endogenous treatment timing;
\item dynamic feedback from outcomes to future treatment.
\end{itemize}

\begin{alertblock}{Core lesson}
More fixed effects mean more conditioning, not automatic identification.
\end{alertblock}
\end{frame}

\section{Random Effects}

\begin{frame}{Random Effects Model}
Random effects also starts from
\[
Y_{it}=\beta D_{it}+X_{it}'\gamma+\alpha_i+u_{it},
\]
but assumes
\[
\Cov(\alpha_i,D_{it})=0
\]
and analogous orthogonality with regressors.

Under the model, GLS can exploit both within- and between-unit variation.

\begin{columns}[T]
\begin{column}{0.48\textwidth}
\begin{block}{Potential advantage}
Greater efficiency when the assumptions are credible.
\end{block}
\end{column}
\begin{column}{0.48\textwidth}
\begin{alertblock}{Potential failure}
Bias when persistent unit heterogeneity correlates with treatment.
\end{alertblock}
\end{column}
\end{columns}
\end{frame}

\begin{frame}{Fixed Effects vs Random Effects}
The difference is not merely computational.

\begin{tabular}{p{3.2cm}p{4.3cm}p{4.3cm}}
\toprule
& \textbf{Fixed effects} & \textbf{Random effects} \\
\midrule
\(\alpha_i\) correlation with regressors & Allowed & Ruled out by model \\
Variation used & Within-unit & Within + between \\
Time-invariant regressors & Not separately estimated & Can be estimated \\
Main risk & Low within variation & Invalid orthogonality \\
\bottomrule
\end{tabular}

\vspace{0.4cm}
Substantive knowledge about treatment assignment should drive the choice before any specification test.
\end{frame}

\begin{frame}{Hausman Comparison Logic}
A classical Hausman-style comparison uses the idea that:
\begin{itemize}
\item FE remains consistent under correlation between \(\alpha_i\) and regressors, under its own assumptions;
\item RE is inconsistent when that correlation is present;
\item a systematic FE--RE coefficient difference may therefore signal RE misspecification.
\end{itemize}

\begin{alertblock}{Do not turn this into a mechanical rule}
Specification tests do not replace a causal argument about the treatment-assignment mechanism.
\end{alertblock}
\end{frame}

\section{Uncertainty in Panel Models}

\begin{frame}{Why IID Standard Errors Are Often Wrong}
Within a unit, disturbances may be serially correlated:
\[
\Cov(u_{it},u_{is})\neq 0
\qquad (t\neq s).
\]

There may also be heteroskedasticity:
\[
\Var(u_{it}\mid X,D)\neq \sigma^2.
\]

Naive IID standard errors can then be too small.

\begin{block}{Common response}
Cluster standard errors at the level where treatment assignment and residual dependence plausibly occur.
\end{block}
\end{frame}

\begin{frame}{Clustered Standard Errors}
For unit-level clustering, arbitrary covariance is allowed within a unit while units are treated as approximately independent across clusters.

\begin{itemize}
\item cluster at the treatment-assignment or dependence level when possible;
\item few clusters require special caution;
\item two-way clustering can be appropriate when dependence runs across units and time;
\item standard-error choice should follow the design, not convention.
\end{itemize}

\begin{alertblock}{Important}
Clustering repairs an uncertainty calculation under dependence. It does not repair endogeneity.
\end{alertblock}
\end{frame}

\begin{frame}{Cross-Sectional Dependence}
Units can share shocks not fully absorbed by time fixed effects:
\[
\Cov(u_{it},u_{jt})\neq 0.
\]

Possible sources include:
\begin{itemize}
\item geographic spillovers;
\item common supply chains;
\item network effects;
\item sector-level shocks;
\item policy coordination.
\end{itemize}

This can affect both standard errors and the validity of the identifying comparison.
\end{frame}

\section{Time-Varying Confounding}

\begin{frame}{Fixed Effects Remove Only Time-Invariant Confounding}
Suppose
\[
Y_{it}=\beta D_{it}+\delta Z_{it}+\alpha_i+u_{it},
\]
and \(Z_{it}\) varies over time.

If \(Z_{it}\) affects both treatment and outcome, unit FE does not remove it.

\[
\Cov(D_{it},Z_{it})\neq 0
\]
can still induce omitted-variable bias.

\begin{alertblock}{Diagnostic question}
Which confounders are stable enough to be absorbed by FE, and which evolve with the treatment process?
\end{alertblock}
\end{frame}

\begin{frame}{Unit-Specific Trends}
Sometimes units would have evolved differently even without treatment.

A richer model may include
\[
Y_{it}=\beta D_{it}+\alpha_i+\lambda_t+\theta_i t+u_{it}.
\]

But unit-specific trends:
\begin{itemize}
\item can absorb meaningful treatment variation;
\item impose a particular functional form;
\item may increase variance substantially;
\item can obscure rather than solve misspecification.
\end{itemize}

\begin{block}{Principle}
Trend controls need substantive justification, not automatic inclusion.
\end{block}
\end{frame}

\section{Dynamic Panels}

\begin{frame}{Lagged Outcomes Change the Problem}
A dynamic panel may contain
\[
Y_{it}=\rho Y_{i,t-1}+\beta D_{it}+\alpha_i+u_{it}.
\]

After the within transformation,
\[
Y_{i,t-1}-\bar Y_{i,-1}
\]
is mechanically correlated with
\[
u_{it}-\bar u_i.
\]

This creates finite-\(T\) bias in the FE estimator, commonly called Nickell bias.
\end{frame}

\begin{frame}{Dynamic Treatment Effects}
Treatment may affect outcomes over several periods:
\[
Y_{it}
=
\alpha_i+\lambda_t+
\sum_{k=0}^{K}\beta_kD_{i,t-k}
+
u_{it}.
\]

Relevant estimands can include:
\begin{itemize}
\item immediate effect \(\beta_0\);
\item delayed effects \(\beta_k\);
\item cumulative effect \(\sum_{k=0}^{K}\beta_k\).
\end{itemize}

A static treatment coefficient can be misleading when effects persist or accumulate.
\end{frame}

\section{Applied Reasoning}

\begin{frame}{Example: Regional Pricing Policy}
Suppose regions are observed monthly and some change pricing policy.

A baseline panel model is
\[
Y_{it}=\beta D_{it}+\alpha_i+\lambda_t+u_{it}.
\]

Before interpreting \(\widehat\beta\) causally, ask:
\begin{itemize}
\item Did policy change because demand was already shifting?
\item Are there region-specific macro shocks?
\item Is treatment measured without substantial error?
\item Does policy spill into neighboring regions?
\item Are effects immediate or lagged?
\end{itemize}
\end{frame}

\begin{frame}{A Panel-Econometrics Workflow}
\begin{enumerate}
\item Describe the panel structure and treatment timing.
\item Plot within-unit and between-unit variation.
\item State which unobserved factors FE is intended to absorb.
\item State the required exogeneity condition.
\item Choose one-way or two-way FE based on the shock structure.
\item Choose clustering based on assignment and dependence.
\item Probe time-varying confounding and trends.
\item Check dynamic feedback and lagged effects.
\item Report the population and variation that identify the coefficient.
\end{enumerate}
\end{frame}

\begin{frame}{Common Failure Modes}
\begin{itemize}
\item interpreting FE as a universal cure for omitted variables;
\item estimating effects from almost no within-unit treatment variation;
\item adding fixed effects without understanding what variation remains;
\item using IID standard errors in strongly dependent panels;
\item selecting RE because it is more efficient despite implausible orthogonality;
\item ignoring dynamic treatment assignment;
\item confusing statistical precision with causal identification.
\end{itemize}
\end{frame}

\begin{frame}{Synthesis: What Fixed Effects Actually Buy}
\begin{itemize}
\item Panel data let us compare units with themselves over time.
\item Unit FE exactly remove time-invariant additive heterogeneity.
\item Two-way FE additionally remove shocks common to all units at each time.
\item Causal interpretation still requires exogeneity of the remaining within-unit treatment variation.
\item Time-varying confounding, anticipation, spillovers, and feedback remain substantive threats.
\item Panel dependence must be reflected in uncertainty estimates.
\item Dynamic panels require methods beyond static FE when lagged outcomes are endogenous after transformation.
\end{itemize}
\end{frame}

\begin{frame}{Selected References}
\small
\begin{itemize}
\item Wooldridge, J. M. (2010). \textit{Econometric Analysis of Cross Section and Panel Data}, 2nd ed. MIT Press.
\item Baltagi, B. H. (2021). \textit{Econometric Analysis of Panel Data}, 6th ed. Springer.
\item Arellano, M. (2003). \textit{Panel Data Econometrics}. Oxford University Press.
\item Angrist, J. D., and Pischke, J.-S. (2009). \textit{Mostly Harmless Econometrics}. Princeton University Press.
\item Nickell, S. (1981). Biases in Dynamic Models with Fixed Effects. \textit{Econometrica}, 49(6), 1417--1426.
\item Cameron, A. C., and Miller, D. L. (2015). A Practitioner's Guide to Cluster-Robust Inference. \textit{Journal of Human Resources}, 50(2), 317--372.
\end{itemize}
\end{frame}

\contactslide

\end{document}
